\SourceLeaf{464}
\par\Needspace{8\baselineskip}\bigskip
\pdfbookmark[3]{Exposition, 15}{pallas-II-article15}
\begin{center}15.\end{center}
\noindent\hspace*{\parindent}The preceding investigations are independent of the nature of the perturbing forces. Suppose now that they arise from the action of another body on both the planet and the Sun. Let \(\mu m\) be the mass of this body, \(x',y',z'\) its coordinates with respect to the three fundamental planes, \mbox{\(r'=\BareRoot{(x'x'+y'y'+z'z')}\)} its distance from the Sun, and \(\rho=\BareRoot{((x'-x)^2+(y-y')^2+(z-z')^2)}\) its distance from the perturbed planet. By known principles, we then have
\[
\begin{aligned}
\xi&=\dfrac{\mu mx'}{r'^3}-\dfrac{\mu m(x'-x)}{\rho^3}\\
\eta&=\dfrac{\mu my'}{r'^3}-\dfrac{\mu m(y'-y)}{\rho^3}\\
\zeta&=\dfrac{\mu mz'}{r'^3}-\dfrac{\mu m(z'-z)}{\rho^3}\mathpunct{\cdot}
\end{aligned}
\]
\par Denoting the heliocentric position of the perturbing planet by \(L'\), we have
\[
\begin{aligned}
x'&=r'\cos XL'\\
y'&=r'\cos YL'\\
z'&=r'\cos ZL';
\end{aligned}
\]
so the formulas of Article 9 give the following, taking into account that \(\Lambda L=ML=90^0\)
\[
\begin{aligned}
T&=\mu\left(\dfrac{1}{\rho^3}-\dfrac{1}{r'^3}\right)r'\cos LL'-\dfrac{\mu r}{\rho^3}\\
V&=\mu\left(\dfrac{1}{\rho^3}-\dfrac{1}{r'^3}\right)r'\cos\Lambda L'\\
W&=\mu\left(\dfrac{1}{\rho^3}-\dfrac{1}{r'^3}\right)r'\cos ML'.
\end{aligned}
\]
Let \(\beta'\) be the inclination of the radius vector \(r'\) to the plane of the osculating orbit of the perturbed planet (taken with the sign \(+\) on the side of the pole \(M\)). Let the projection of the radius vector onto this plane form, with the line of this plane's ascending node in the \(x,y\) coordinate plane, the angle \(v'-\AscendingNode\), reckoned in the direction of motion of the perturbed planet. Then \(\beta',v'\) can be regarded as the heliocentric latitude and longitude of the perturbing body relative \SourceLeaf{465}to the plane of the osculating orbit of the perturbed planet. It is easy to see that
\[
\begin{aligned}
\cos LL'&=\cos\beta'\cos(v'-v)\\
\cos\Lambda L'&=\cos\beta'\sin(v'-v)\\
\cos ML'&=\sin\beta'.
\end{aligned}
\]
\par Finally, to gather here everything needed for the calculation, suppose that the plane in which the perturbing body moves is inclined to the \(x,y\) coordinate plane by \(i'\), and let \(\AscendingNode'\) be the longitude of the ascending node measured from the \(x\) axis. Let \(w'-\AscendingNode'\) be the distance of the heliocentric position \(L'\) from that same node, so that \(w'\) is what is called its true longitude in the orbit. Determine \(N,\Delta,J\) from the following equations.
\[
\begin{aligned}
\sin\tfrac12J\mathbin{.}\sin\tfrac12(\Delta+N)&=\sin\tfrac12(i'+i)\mathbin{.}\sin\tfrac12(\AscendingNode'-\AscendingNode)\\
\sin\tfrac12J\mathbin{.}\cos\tfrac12(\Delta+N)&=\sin\tfrac12(i'-i)\mathbin{.}\cos\tfrac12(\AscendingNode'-\AscendingNode)\\
\cos\tfrac12J\mathbin{.}\sin\tfrac12(\Delta-N)&=\cos\tfrac12(i'+i)\mathbin{.}\sin\tfrac12(\AscendingNode'-\AscendingNode)\\
\cos\tfrac12J\mathbin{.}\cos\tfrac12(\Delta-N)&=\cos\tfrac12(i'-i)\mathbin{.}\cos\tfrac12(\AscendingNode'-\AscendingNode).
\end{aligned}
\]
\par This done, we have
\[
\begin{aligned}
\operatorname{tang}(v'-\AscendingNode-\Delta)&=\cos J\mathbin{.}\operatorname{tang}(w'-\AscendingNode'-N)\\
\operatorname{tang}\beta'&=\operatorname{tang}J\mathbin{.}\sin(v'-\AscendingNode-\Delta)\\
\sin\beta'&=\sin J\mathbin{.}\sin(w'-\AscendingNode'-N)\\
\rho&=\BareRoot{(rr+r'r'-2rr'\cos\beta'\cos(v'-v))}.
\end{aligned}
\]
